\section{Synthesis by a nonnegative charge density}\label{sec:density}

We now construct an attractive Coulomb potential whose low-energy
matrix realizes the modes and hopping of Section~\ref{sec:wells}.
A broad slab supplies the transverse oscillator near the modes and a
positive background density. Adding the planar wells then changes the
density by a bounded, compactly supported term, so the total density
remains nonnegative. We first construct this potential, then establish
a fixed energy cost outside the mode space and compute its matrix on
that space.

\subsection{The slab and the well-depth counterterm}

Use integers \(S\) of size between \(D\) and \(10D\), and \(H=\lceil\lambda^5\rceil\). (For \(S\) one can use a binary bit-length scale.) Let
\[
 \Omega=[-H,H]^2\times[-S,S],\qquad
 F_H(y)=-\rho\int_\Omega\frac{db}{|y-b|},\qquad C_H=F_H(0).
\]
The occupancy identity~\eqref{eq:occupancy-completion} requires the
one-body diagonal \(-s_i/\lambda\), where
\(s_i=\sum_{j\ne i}v_{ij}\). Put \(W_i=W(\,\cdot-u_i)\).
Since \(\langle\psi_i,W_i\psi_i\rangle=\int\varphi^2W<0\),
a small increase in each well depth has the required sign and size.
Take a fixed sufficiently smooth cutoff \(\zeta\) between 0 and 1,
supported inside \((-1,1)\), equal to 1 for \(|z|\le 1/2\);
it may be piecewise polynomial with rational data. Define
\[
 b_i=\frac{s_i}{-\int\varphi^2 W},\qquad
 V_\ell(r,z)=\zeta(z/S)\sum_i(1+b_i/\lambda)W_i(r).
\]
These \(b_i\ge0\) are polynomially bounded independent of \(D\); hence \(0\le b_i/\lambda<1\) at our scales. For all \(S\ge1\) the derivatives we need of \(V_\ell\) are bounded independently, by disjoint planar supports. Choose the fixed integer \(\rho\) so large, using \(1+b_i/\lambda\le2\), that \(|\Delta V_\ell/(4\pi)|\le\rho/2\). There is no circularity, since the bound based on \(2,W,\zeta\) does not depend on the vertical Gaussian or \(b_i\). The centers fit well inside the horizontal box for our scale. Thus
\[
               F_H+V_\ell
   =-\int_\Omega \frac{\rho+\Delta V_\ell(b)/(4\pi)}{|y-b|}\,db       \tag{4}\label{eq:4}
\]
by compact support and the fundamental solution identity.

Thus the density in~\eqref{eq:4} lies between \(\rho/2\) and
\(3\rho/2\) on \(\Omega\). Its total mass is \(\rho|\Omega|\),
because \(\int\Delta V_\ell=0\). The next estimates show that the
truncated slab approximates the harmonic potential near the modes and
provides a global lower bound outside that region.

Here are useful background estimates with constants independent of \(n,\lambda\):
\[
 \begin{split}
 F_H(y)-C_H &\ge 2\pi\rho \min(z^2,S^2)- C S^3/H,\\
 |F_H(y)-C_H-2\pi\rho z^2|
      &\le C(S^3+S|r|^2)/H
           \qquad (|r|\le H/2,\ |z|\le S).
 \end{split}                                                       \tag{5}\label{eq:5}
\]
Indeed \(F_H\) at each \(z\) is minimized at \(r=0\), and is also nondecreasing in \(|z|\) for fixed \(r\). These statements follow one coordinate at a time by integration of even decreasing kernels on centered intervals. At \(r=0\), extending the horizontal integration to the plane in the difference between \(z\) and 0 gives for \(|z|\le S\) exactly \(2\pi\rho\int_{-S}^S(|z-w|-|w|)\,dw=2\pi\rho z^2\), by polar integration. The omitted difference costs \(CS^3/H\), by inverse cube decay outside horizontal radius \(H\). For \(|r|\le H/2\), the change relative to the horizontal center is bounded by \(CS|r|^2/H\). For example translate the integration square instead of translating the kernel. The first derivatives then come from sides at distance at least \(H/2\), and second derivatives are bounded by \(CS/H\) by integrating first derivatives of the kernel on the sides. Evenness at zero gives the bound. Finally \(0\ge F_H(y)\ge C_H\) and \(|C_H|\le CSH\).

\subsection{The one-particle gap and residual}

The potential has now been expressed as the Coulomb field of a
nonnegative density. Before replacing it by point nuclei we need two
properties: a fixed energy gap outside the prescribed modes, and a small
residual on those modes. The gap will control arbitrary continuum
states, while the residual will control mixing with the mode space.

Let the shifted continuous one-particle operator, still before point-charge replacement, be
\[
       L_c=-\Delta_y/2+F_H+V_\ell-C_H-e_* .
\]
We suppress the tensor identity on spin. Let \(P\) project onto the span
of the \(\psi_i\), including both spins when present, and write
\(\Gamma_{ij}=\langle\psi_i,\psi_j\rangle\) for their spatial Gram
matrix. Exponential tails and the \(0.95D\) spacing give
\(\|\Gamma-I\|\le p(n,D)\lambda^{-0.8}\): for example, each
off-diagonal overlap is \(O(e^{-0.9|u_i-u_j|})\), by leaving part of
the decay as an integrable factor. We use the orthonormal modes
\[
 \widehat\psi_j=\sum_i\psi_i(\Gamma^{-1/2})_{ij},\qquad
 \|\Gamma^{-1/2}-I\|\le p(n,D)\lambda^{-0.8}.
\]
In particular, unit vectors in \(\operatorname{Ran}P\) have
polynomially bounded one-particle kinetic energy.

\begin{proposition}[Continuous one-particle control]\label{prop:continuous-control}
For the density and separated centers constructed above, with the
polynomial scale chosen sufficiently large, there are fixed constants
\(c>0,C<\infty\) such that
\[
 \begin{split}
 (I-P)L_c(I-P)&\ge c(I-P)\quad\text{on the complement},\\
 L_c&\ge -\Delta_y/2-C,\qquad
 \|L_cP\|\le p(n,D)\lambda^{-0.9}.
 \end{split}                                                       \tag{6}\label{eq:6}
\]
The inequalities are quadratic-form inequalities on the one-particle
$H^1$ domain, and the last bound is an operator norm from the range of
$P$ to $L^2$. All estimates hold with either spin included.
\end{proposition}

\begin{proof}
For the first inequality compare below, using \eqref{eq:5}, with the separable planar multiwell \(-\Delta_r/2+\sum W_i\) plus the vertical operator with capped potential \(2\pi\rho\min(z^2,S^2)\), shifted by \(-e_*\). Errors in the lower bound are \(O(S^3/H)+O(p(n,D)/\lambda)\), since \(\zeta W_i\ge W_i\) and the \(b_i/\lambda\) terms are small and bounded everywhere. 

For clarity we detail the gap argument. In the plane use a quadratic partition of unity consisting of cutoffs about each site (equal to one inside radius \(D/8\), supported at radius \(D/3\)) and a remaining exterior function. They can be built by smooth angle interpolation from 0 to \(\pi/2\) on the transition annuli using sine and cosine, with squared gradients summing to \(C/D^2\) or less. The localization form identity (IMS localization; compare
\cite[Lemma 3.1]{Simon1983}), obtained here by the product rule with
kinetic coefficient $1/2$, then costs \(O(1/D^2)\). In each site term use the isolated ground and gap bound, and in the exterior use the free kinetic lower bound. This gives
\[
 -\Delta_r/2+\sum W_i \ \ge\
 -1/2+c_1(I-P_r)-o(1)I,
\]
where \(P_r\) projects on the planar translates of \(\varphi\), \(c_1>0\) can be taken below the isolated gap and \(1/2\), and the small bound can be made uniformly small. In detail, with site cutoffs \(\theta_i\) and \(\varphi_i=\varphi(\cdot-u_i)\), the site term evaluated on \(\theta_i u\) is at least
\[
 (-1/2+c_1)\|\theta_i u\|^2-c_1|\langle\theta_i\varphi_i,u\rangle|^2;
\]
the exterior term also bounds its share of the \(-1/2+c_1\) term as there is no well on its support. The sum of rank-one operators from the cutoff ground-state modes differs from \(P_r\) by at most \(O(m e^{-c' D})+O(p(n,D)\lambda^{-0.8})\), by the tails and Gram bound. 

Similarly localize the vertical problem with an inner cutoff equal to one on \(|z|\le S/2\), supported on \(|z|\le S\), and an exterior function. Use the harmonic gap and the large \(z\) potential lower bound on the respective terms, at gradient cost \(O(S^{-2})\), with Gaussian loss on the rank-one mode. This gives a vertical lower bound \(\omega/2+c_2(I-P_g)-o(1)I\), where \(P_g=|g\rangle\langle g|\). Since \(P=P_r\otimes P_g\) on the spatial space and
\[
 (I-P_r)\otimes I+I\otimes(I-P_g)\ge I-P,
\]
the tensor sum has a fixed gap on \(I-P\) once the errors are small.

The slowly small errors such as \(O(D^{-2})\) here are needed only for a fixed gap estimate, not as the energy accuracy per particle.

The kinetic lower bound in \eqref{eq:6} follows from \eqref{eq:5} or background minimality and bounded \(V_\ell\). For the residual estimate, \eqref{eq:5}, the crude global background bound, and the mode tails give
\[
  \|(F_H-C_H-2\pi\rho z^2)\psi_i\|_2\le p(n,D)\lambda^{-2}.
\]
Indeed moments with \(|r|\le H/2,\ |z|\le S\) suffice there, and outside this region the Gaussian or planar exponential beats the global polynomial magnitudes, using \(S\ge\log\lambda\). The cutoff defect in \(V_\ell\) versus \(\sum_i(1+b_i/\lambda)W_i\) acting on any mode is likewise negligible. Since \((-{\Delta_y}/2+2\pi\rho z^2+W_j-e_*)\psi_j=0\), the remaining residual is a sum of other wells acting on the small planar tails and the term with \(b_j/\lambda\). This proves \eqref{eq:6}.
\end{proof}

\subsection{The compressed one-particle matrix}

The fixed complementary gap alone does not give the required energy
accuracy. We therefore compute the matrix on the prescribed modes,
including its diagonal counterterm, to a sharper scale.

The continuous one-body compression can now be evaluated more precisely. On the nonorthogonal functions before Gram correction, the column on mode \(j\) comes, up to background and cutoff errors as above, from
\[
                  \Big(\sum_{k\ne j} W_k+\lambda^{-1}\sum_k b_k W_k\Big)\psi_j .
\]


On row \(j\), the leading contribution is \(-\lambda^{-1}\sum_{i\ne j}v_{ij}\) by construction; the other-well integrals there use two tail factors. On row \(i\ne j\) the first sum at \(k=i\) gives \(-T(|u_i-u_j|)\). The cases \(k\ne i,j\) of the first sum are products of two tail factors at a third site; each such integral is \(O(\lambda^{-1.8})\) by the spacing and compact well support. Terms with the explicit \(\lambda^{-1}\) give only negligible off-diagonal errors (at least one planar tail on each well support, hence bounded by \(p(n,D)\lambda^{-1.9}\) altogether on such an entry). Nonlink hoppings cost at most \(p(n,D)\lambda^{-1.2}\) per term. 

To pass to the orthonormal modes, let
\(A_{ij}=\langle\psi_i,L_c\psi_j\rangle\). The matrix becomes
\(\Gamma^{-1/2}A\Gamma^{-1/2}\), and
\[
 \|\Gamma^{-1/2}A\Gamma^{-1/2}-A\|
 \le C\|\Gamma-I\|\,\|A\|
 \le p(n,D)\lambda^{-1.7}
\]
by~\eqref{eq:6}. Here \(A\) is the \emph{shifted} matrix: subtracting
\(C_H+e_*\) before making this estimate removes the large background
scalar. The correction therefore remains negligible after
multiplication by \(\lambda\).

Define the spatial matrix
\[
 \mathsf M_{ij}=\begin{cases}
  -s_i,&i=j,\\
  -\lambda T(|u_i-u_j|),&\{i,j\}\text{ is a link},\\
  0,&\text{otherwise}.
 \end{cases}
\]
The preceding estimates give
\[
 \bigl\|\lambda\Gamma^{-1/2}A\Gamma^{-1/2}-\mathsf M\bigr\|
       \le p(n,D)\lambda^{-0.1}.
\]
This is the required one-particle compression, with the identity on
spin. The diagonal supplies the depth counterterm, and the link
entries supply the calibrated hopping.
